Execution feedback can reveal that a program is wrong without identifying why it is wrong. A single failed test may support incompatible explanations, yet a repair generator must ultimately choose a coherent edit. We propose Particle-Belief Program Filtering (PBPF), a framework that separates uncertainty about a candidate's failure mechanism from uncertainty over candidate programs. PBPF maintains a candidate-specific particle posterior over a learned latent failure state. Ordered observations update that state through a calibrated categorical likelihood while the code remains fixed; a new program version alone triggers a transition and an importance-corrected proposal. A sampled posterior particle is then mapped to a compact soft prefix and held fixed throughout one complete repair from a frozen code actor. The method is not training-free: belief learning and prefix learning are separate stages. We specify a prospective evaluation that first tests posterior validity and withheld-outcome prediction, then tests hidden-test repair under shared candidate banks and equal generation budgets. Matched deterministic memories, evidence interventions, and coherence ablations are designed to separate useful diagnosis from additional computation. The package contains no completed confirmatory results; predictive gains, repair gains, calibration, and replication outcomes remain explicitly pending. The intended claim is a testable link between sequential diagnosis and coherent repair, not priority for Bayesian debugging.
